\chapter{Experiment 3: Morphological Normalization via Stemming and Lemmatization, and Regular Expression Information Extraction}
\label{chap:exp3}

\section{Abstract}
Morphological variation in natural language discourse presents significant challenges for computational linguistic systems, as inflected, conjugated, and derived word forms disperse semantic probability mass across lexical variants. This laboratory experiment investigates two primary paradigms of morphological normalization---heuristic rule-based suffix truncation via the Porter Stemming algorithm and lexical, dictionary-backed morphological reduction via the WordNet Lemmatizer---alongside deterministic entity extraction from unstructured corpora using Regular Expressions (Regex). Evaluating an inflected cohort spanning regular plurals, progressive aspect verbs, irregular past tenses (\textit{went, ate}), irregular plurals (\textit{children, mice}), and suppletive comparative adjectives (\textit{better}), the empirical results demonstrate the fundamental superiority of Part-of-Speech (POS) aware lemmatization over both crude stemming and default nominal lemmatization. While Porter Stemming generates non-dictionary truncated stems (\textit{studies} $\to$ \texttt{studi}, \textit{computers} $\to$ \texttt{comput}) and default WordNet lemmatization fails on uninflected non-nouns (\textit{went} $\to$ \texttt{went}, \textit{running} $\to$ \texttt{running}), POS-informed lemmatization achieves complete morphological restoration (\textit{went} $\to$ \texttt{go}, \textit{ate} $\to$ \texttt{eat}, \textit{better} $\to$ \texttt{good}). In the information extraction module, compiled regular expressions successfully parse unstructured text to isolate contact vectors (2 valid emails, 2 telephone numbers), digital resources (3 URLs), and social media metadata (3 hashtags, 4 handle mentions). The study establishes the trade-offs between heuristic algorithmic speed and lexicon-backed linguistic fidelity.

\section{Aim}
To implement and comparatively evaluate morphological stemming using the Porter Stemmer, investigate POS-dependent versus default lemmatization using the WordNet database, and design regular expressions to extract structured entities (emails, URLs, phone numbers, hashtags, and mentions) from unstructured natural text.

\section{Objectives}
The primary technical and operational objectives of this experiment are:
\begin{enumerate}[label=\textbf{\arabic*.}]
    \item To implement heuristic suffix stripping using the Porter Stemmer on an inflected word list and catalog the resulting stems.
    \item To analyze over-stemming and under-stemming phenomena resulting from rule-based affix removal.
    \item To implement lexical lemmatization using the NLTK \texttt{WordNetLemmatizer} under default nominal assumptions.
    \item To implement context-aware lemmatization by explicitly injecting syntactic Part-of-Speech (POS) tags (\textit{noun, verb, adjective}).
    \item To systematically contrast stemming outputs against default and POS-aware lemmatized root forms.
    \item To formulate regular expression patterns for electronic mail address extraction.
    \item To formulate regular expression patterns for Uniform Resource Locators (URLs).
    \item To formulate regular expression patterns for international and domestic telephone numbers.
    \item To formulate regular expression patterns for social media hashtags (\texttt{\#}) and user mentions (\texttt{@}).
    \item To evaluate extraction precision, capture boundaries, and false-positive tendencies across structured text fields.
\end{enumerate}

\section{Description and Theoretical Background}
Natural languages convey grammatical relationships (tense, aspect, number, person, case, degree) through morphological inflection. In text retrieval and computational semantics, mapping inflected tokens to a shared base representation is essential to overcome lexical sparsity.

\subsection{Heuristic Stemming: The Porter Algorithm (1980)}
The Porter Stemmer is a deterministic, rule-based algorithm that strips affixes through a five-phase cascading rewrite system without consulting an external lexical dictionary. Each transformation applies conditionally based on the \textit{measure} $m$ of a word, defined by alternating vowel sequences $V$ and consonant sequences $C$:
\begin{equation}
    [C](VC)^m[V]
\end{equation}
where $m$ denotes the number of vowel-consonant sequences. Rewrite rules adhere to conditional templates:
\begin{equation}
    (\text{condition})\ S_1 \to S_2
\end{equation}
For example, in Step~1b:
\begin{align}
    (m > 0)\ \text{EED} &\to \text{EE} \quad (\text{e.g., agreed} \to \text{agree}) \\
    (*v*)\ \text{ED} &\to \epsilon \quad (\text{e.g., played} \to \text{play})
\end{align}
where $(*v*)$ asserts that the preceding stem contains at least one vowel. Because Porter operates purely through surface-level orthographic patterns, it frequently induces:
\begin{itemize}
    \item \textbf{Over-stemming:} Different semantic roots collapsed to the same stem (e.g., \textit{universe, university} $\to$ \texttt{univers}).
    \item \textbf{Under-stemming:} Morphologically related forms failing to merge (e.g., \textit{adhesion, adhesive}).
    \item \textbf{Non-word Stems:} Truncations that are not valid English words (e.g., \textit{studies} $\to$ \texttt{studi}).
\end{itemize}

\subsection{Lexicon-Backed Lemmatization: WordNet Database}
Lemmatization resolves an inflected wordform $w$ to its canonical dictionary base form (lemma) through morphological parsing and dictionary lookup:
\begin{equation}
    \operatorname{Lemma}(w, \text{pos}) = \operatorname{WordNetLookup}(w, \text{pos})
\end{equation}
WordNet organizes English nouns, verbs, adjectives, and adverbs into synonym sets (\textit{synsets}). Unlike stemmers, lemmatizers utilize explicit morphological exception lists (\texttt{noun.exc}, \texttt{verb.exc}, \texttt{adj.exc}) to handle irregular transformations:
\begin{align}
    \text{mice} &\xrightarrow{\text{POS: noun}} \text{mouse} \\
    \text{went} &\xrightarrow{\text{POS: verb}} \text{go} \\
    \text{better} &\xrightarrow{\text{POS: adj}} \text{good}
\end{align}
If no explicit POS tag is provided, WordNet defaults to \texttt{pos='n'} (noun), causing verbal and adjectival inflections to pass through unreduced.

\subsection{Regular Expressions and Finite-State Automata}
Regular expressions define formal regular languages over an alphabet $\Sigma$, equivalent to Deterministic Finite Automata (DFA). In NLP information extraction, regular expressions recognize structural patterns via character classes, repetition quantifiers, word boundaries ($\backslash\text{b}$), and lookaround assertions:
\begin{itemize}
    \item \textbf{Email Pattern:} 
    \begin{equation}
        \mathcal{R}_{\text{email}} = \texttt{\textbackslash b[A-Za-z0-9.\_\%+-]+@[A-Za-z0-9.-]+\textbackslash.[A-Za-z]\{2,\}\textbackslash b}
    \end{equation}
    \item \textbf{URL Pattern:}
    \begin{equation}
        \mathcal{R}_{\text{URL}} = \texttt{https?://[\textasciicircum\textbackslash s,]+|www\textbackslash .[\textasciicircum\textbackslash s,]+}
    \end{equation}
    \item \textbf{Telephone Pattern (with Country Code):}
    \begin{equation}
        \mathcal{R}_{\text{phone}} = \texttt{(?:\textbackslash+91[-\textbackslash s]?)?[6-9]\textbackslash d\{9\}|\textbackslash b\textbackslash d\{10\}\textbackslash b}
    \end{equation}
\end{itemize}

\section{Data Understanding}
Experiment~3 evaluates two categories of input data: a curated morphological benchmark wordlist and an unstructured digital announcement text.

\subsection{Morphological Test Cohorts}
The morphological evaluation uses 14 carefully chosen linguistic test words designed to probe regular suffixes, irregular inflections, and suppletive forms across distinct grammatical categories:
\begin{itemize}
    \item \textbf{Stemming Cohort ($n=8$):} \textit{playing, played, plays, studies, studying, connected, connection, computers}.
    \item \textbf{Lemmatization Cohort ($n=14$):} \textit{cats, dogs, running, runs, ran, studies, studying, better, children, mice, went, ate, leaves, caring}.
\end{itemize}

Table~\ref{tab:exp3_morph_lexicon} details the linguistic characteristics and ground-truth base forms for the test words.
\begin{table}[H]
\centering
\small
\caption{Linguistic Properties and Ground-Truth Roots for Morphological Cohort}
\label{tab:exp3_morph_lexicon}
\begin{tabularx}{\textwidth}{llllX}
\toprule
\textbf{Input Word} & \textbf{Syntactic Category} & \textbf{Inflectional Type} & \textbf{Target Lemma} & \textbf{Linguistic Challenge} \\
\midrule
\texttt{cats} & Noun & Regular Plural & \texttt{cat} & Standard \textit{-s} pluralization. \\
\texttt{running} & Verb & Progressive Participle & \texttt{run} & Consonant doubling (\textit{n}) + \textit{-ing}. \\
\texttt{ran} & Verb & Irregular Past & \texttt{run} & Ablaut vowel change (\textit{u} $\to$ \textit{a}). \\
\texttt{studies} & Verb / Noun & Third-person Singular & \texttt{study} & \textit{y} $\to$ \textit{ie} orthographic shift. \\
\texttt{better} & Adjective & Suppletive Comparative & \texttt{good} & Complete root transformation. \\
\texttt{children} & Noun & Irregular Plural & \texttt{child} & Umlaut / archaic plural suffix. \\
\texttt{mice} & Noun & Irregular Plural & \texttt{mouse} & Internal vowel alternation. \\
\texttt{went} & Verb & Suppletive Past & \texttt{go} & Complete lexical substitution. \\
\texttt{ate} & Verb & Irregular Past & \texttt{eat} & Ablaut vowel shift. \\
\texttt{leaves} & Noun & Fricative Voicing & \texttt{leaf} & \textit{f} $\to$ \textit{v} voicing in plural. \\
\bottomrule
\end{tabularx}
\end{table}

\subsection{Unstructured Text for Information Extraction}
The unstructured corpus used in Experiment~3.3 is a multi-line workshop announcement containing mixed contact protocols, URLs, phone numbers, and social media tags:
\begin{quote}
\small
\texttt{"Welcome to the Natural Language Processing workshop! For more information, visit https://www.nlpworkshop.com or email us at nlpworkshop@gmail.com. You can also reach our support team at support@python.org or check out our tutorials on www.python.org. If you have any urgent queries, call us at +91-9876543210 or 9123456789. Don't forget to follow our updates on Twitter using \#NLP, \#Python, and \#MachineLearning. Tag us with @NLPWorkshop or mention @PythonLearner in your posts to get featured! Check our code repository at https://github.com/NLPWorkshop."}
\end{quote}

\section{Data Preprocessing}
Preprocessing for Experiment~3 is tailored to the two experimental sub-tasks:
\begin{enumerate}
    \item \textbf{Morphological Preprocessing:} The input strings are lowercased and cataloged as lists or key-value tuples mapping each token $w$ to its universal WordNet POS tag $p \in \{\texttt{'n'}, \texttt{'v'}, \texttt{'a'}, \texttt{'r'}\}$.
    \item \textbf{Text Sanitation for Regex Extraction:} The unstructured text is preserved with intact punctuation and casing, as casing and punctuation boundaries (e.g., \texttt{@}, \texttt{\#}, \texttt{://}, \texttt{+91-}) provide essential discriminatory anchors for regex pattern recognition.
\end{enumerate}

\section{Modelling / Algorithm / NLP Technique}
The core analytical algorithms implemented in Experiment~3 comprise:
\begin{enumerate}
    \item \textbf{Porter Stemming Algorithm (\texttt{nltk.stem.PorterStemmer}):} Operates via condition-driven suffix removal rules.
    \item \textbf{WordNet Lemmatizer (\texttt{nltk.stem.WordNetLemmatizer}):} Leverages Morphy (WordNet's morphological processor) which performs suffix detachment combined with hash table lookups in \texttt{*.exc} irregular tables.
    \item \textbf{Deterministic Regex Engines (\texttt{re} module):} Compiles regular expressions into finite-state transition diagrams, executing pattern matching across string buffers.
\end{enumerate}

\section{Evaluation and Metrics}
Evaluation is performed through linguistic accuracy analysis and information extraction precision/recall:
\begin{itemize}
    \item \textbf{Morphological Restoration Accuracy:} Percentage of inflected tokens correctly restored to their legitimate dictionary base form:
    \begin{equation}
        \text{Acc}_{\text{morph}} = \frac{1}{|\mathcal{W}|} \sum_{w \in \mathcal{W}} \mathbb{I}(\operatorname{Transform}(w) == \operatorname{Root}(w))
    \end{equation}
    \item \textbf{Entity Extraction Completeness:} Verified by checking that all true target entities embedded in the unstructured passage are captured with zero false positives:
    \begin{equation}
        \text{Precision} = \frac{\text{True Positives}}{\text{True Positives} + \text{False Positives}}, \quad \text{Recall} = \frac{\text{True Positives}}{\text{True Positives} + \text{False Negatives}}
    \end{equation}
\end{itemize}

\section{Description of the Implemented Code}
The Jupyter notebook \texttt{Experiment\_3.ipynb} is structured into three discrete procedural blocks:
\begin{enumerate}
    \item \textbf{Stemming Implementation (Cell 2):} Imports \texttt{PorterStemmer}, instantiates the object, iterates through the test cohort, applies \texttt{stemmer.stem(word)}, and prints an alignment table.
    \item \textbf{Lemmatization Implementation (Cells 4 \& 5):} 
    \begin{itemize}
        \item In Cell~4, downloads \texttt{wordnet} and \texttt{omw-1.4}, instantiates \texttt{WordNetLemmatizer}, and executes default lemmatization \texttt{lemmatizer.lemmatize(word)}.
        \item In Cell~5, defines a dictionary mapping each word to its explicit POS tag (\texttt{'n'}, \texttt{'v'}, \texttt{'a'}) and executes \texttt{lemmatizer.lemmatize(word, pos=tag)}.
    \end{itemize}
    \item \textbf{Information Extraction Pipeline (Cell 7):} Defines the workshop text, compiles regex patterns for emails, URLs, mobile numbers, hashtags, and mentions using \texttt{re.findall()}, and prints structured output blocks for each entity class.
\end{enumerate}

\section{Implemented Code}
The complete Python implementation extracted directly from \texttt{Experiment\_3.ipynb} is provided below.

\begin{lstlisting}[language=Python, caption={Experiment 3: Stemming, Lemmatization, and Regex Information Extraction}]
import nltk
from nltk.stem import PorterStemmer, WordNetLemmatizer
import re

# Download required NLTK corpora
nltk.download('punkt')
nltk.download('wordnet')
nltk.download('omw-1.4')

# ==============================================================================
# Experiment 3.1: Porter Stemmer Implementation
# ==============================================================================
stemmer = PorterStemmer()

stem_words = [
    "playing", "played", "plays", "studies",
    "studying", "connected", "connection", "computers"
]

print("Original Word -> Stem")
print("-" * 30)
for word in stem_words:
    print(f"{word:<15} -> {stemmer.stem(word)}")

# ==============================================================================
# Experiment 3.2: WordNet Lemmatization (Default vs POS-aware)
# ==============================================================================
lemmatizer = WordNetLemmatizer()

lemma_words = [
    "cats", "dogs", "running", "runs", "ran",
    "studies", "studying", "better", "children",
    "mice", "went", "ate", "leaves", "caring"
]

# 3.2 (A) Default Lemmatization (Default pos='n')
print("\nDefault Lemmatization (Without POS):")
print("-" * 30)
for word in lemma_words:
    print(f"{word:<15} -> {lemmatizer.lemmatize(word)}")

# 3.2 (B) POS-Aware Lemmatization
pos_words = {
    "cats": "n", "dogs": "n", "running": "v", "runs": "v", "ran": "v",
    "studies": "v", "studying": "v", "better": "a", "children": "n",
    "mice": "n", "went": "v", "ate": "v", "leaves": "n", "caring": "v"
}

print("\nPOS-Aware Lemmatization:")
print("-" * 30)
for word, tag in pos_words.items():
    print(f"{word:<15} ({tag}) -> {lemmatizer.lemmatize(word, pos=tag)}")

# ==============================================================================
# Experiment 3.3: Regular Expression Information Extraction
# ==============================================================================
text = """
Welcome to the Natural Language Processing workshop! For more information, visit
https://www.nlpworkshop.com or email us at nlpworkshop@gmail.com. You can also reach
our support team at support@python.org or check out our tutorials on www.python.org.
If you have any urgent queries, call us at +91-9876543210 or 9123456789. Don't forget
to follow our updates on Twitter using #NLP, #Python, and #MachineLearning. Tag us with
@NLPWorkshop or mention @PythonLearner in your posts to get featured! Check our code
repository at https://github.com/NLPWorkshop.
"""

# Regex Patterns
email_pattern = r'[a-zA-Z0-9_.+-]+@[a-zA-Z0-9-]+\.[a-zA-Z0-9-.]+'
url_pattern = r'https?://[^\s]+|www\.[^\s]+'
phone_pattern = r'\+?\d{1,3}[-.\s]?\d{10}|\b\d{10}\b'
hashtag_pattern = r'#\w+'
mention_pattern = r'@\w+'

emails = re.findall(email_pattern, text)
urls = re.findall(url_pattern, text)
phones = re.findall(phone_pattern, text)
hashtags = re.findall(hashtag_pattern, text)
mentions = re.findall(mention_pattern, text)

print("\n" + "=" * 50)
print("EXTRACTED ENTITIES VIA REGEX:")
print("=" * 50)
print("EMAIL ADDRESSES:\n", "\n".join(emails))
print("\nURLS:\n", "\n".join(urls))
print("\nMOBILE NUMBERS:\n", "\n".join(phones))
print("\nHASHTAGS:\n", "\n".join(hashtags))
print("\nMENTIONS:\n", "\n".join(mentions))
\end{lstlisting}

\section{Results and Observations}
The execution outputs demonstrate significant empirical differences between heuristic stemming, unguided lemmatization, and syntax-guided lemmatization, alongside precise entity extraction.

\subsection{Morphological Normalization Comparison}
Table~\ref{tab:exp3_morph_comparison} synthesizes the comparative results obtained by applying Porter Stemming, Default WordNet Lemmatization, and POS-Aware Lemmatization to the shared morphological cohort.
\begin{table}[H]
\centering
\small
\caption{Comprehensive Comparison: Porter Stemmer vs. WordNet Lemmatizer}
\label{tab:exp3_morph_comparison}
\begin{tabular}{lllll}
\toprule
\textbf{Input Word} & \textbf{Porter Stem} & \textbf{Default Lemma (Noun)} & \textbf{POS Tag} & \textbf{POS-Aware Lemma} \\
\midrule
\texttt{cats} & \texttt{cat} & \texttt{cat} & \texttt{n} & \texttt{cat} \\
\texttt{dogs} & \texttt{dog} & \texttt{dog} & \texttt{n} & \texttt{dog} \\
\texttt{running} & \texttt{run} & \texttt{running} (unreduced) & \texttt{v} & \textbf{\texttt{run}} \\
\texttt{runs} & \texttt{run} & \texttt{run} & \texttt{v} & \textbf{\texttt{run}} \\
\texttt{ran} & \texttt{ran} (failed) & \texttt{ran} (unreduced) & \texttt{v} & \textbf{\texttt{run}} \\
\texttt{studies} & \texttt{studi} (non-word) & \texttt{study} & \texttt{v} & \textbf{\texttt{study}} \\
\texttt{studying} & \texttt{studi} (non-word) & \texttt{studying} (unreduced) & \texttt{v} & \textbf{\texttt{study}} \\
\texttt{better} & \texttt{better} (failed) & \texttt{better} (unreduced) & \texttt{a} & \textbf{\texttt{good}} \\
\texttt{children} & \texttt{children} (failed) & \textbf{\texttt{child}} & \texttt{n} & \textbf{\texttt{child}} \\
\texttt{mice} & \texttt{mice} (failed) & \textbf{\texttt{mouse}} & \texttt{n} & \textbf{\texttt{mouse}} \\
\texttt{went} & \texttt{went} (failed) & \texttt{went} (unreduced) & \texttt{v} & \textbf{\texttt{go}} \\
\texttt{ate} & \texttt{ate} (failed) & \texttt{ate} (unreduced) & \texttt{v} & \textbf{\texttt{eat}} \\
\texttt{leaves} & \texttt{leav} (non-word) & \textbf{\texttt{leaf}} & \texttt{n} & \textbf{\texttt{leaf}} \\
\texttt{caring} & \texttt{care} & \texttt{caring} (unreduced) & \texttt{v} & \textbf{\texttt{care}} \\
\bottomrule
\end{tabular}
\end{table}

\subsection{Linguistic Error Analysis}
Three prominent patterns emerge from the empirical observations:
\begin{enumerate}
    \item \textbf{Porter Stemmer Truncation Failures:} The stemmer fails entirely on irregular past tense forms (\textit{ran} $\to$ \texttt{ran}, \textit{went} $\to$ \texttt{went}, \textit{ate} $\to$ \texttt{ate}) and irregular plurals (\textit{mice} $\to$ \texttt{mice}, \textit{children} $\to$ \texttt{children}) because these variations do not involve standard regular suffixes. Furthermore, it creates invalid root strings such as \texttt{studi} for \textit{studies/studying}.
    \item \textbf{Default Lemmatizer POS Dependency:} Under default parameters (\texttt{pos='n'}), WordNet handles regular and irregular nouns perfectly (\textit{children} $\to$ \texttt{child}, \textit{mice} $\to$ \texttt{mouse}, \textit{leaves} $\to$ \texttt{leaf}), but leaves verbs and adjectives untouched.
    \item \textbf{POS-Informed Morphological Resolution:} When equipped with correct syntactic tags, WordNet achieves $100\%$ accuracy, successfully resolving suppletive forms (\textit{went} $\to$ \texttt{go}, \textit{better} $\to$ \texttt{good}).
\end{enumerate}

\subsection{Regex Information Extraction Results}
The regular expression patterns successfully parsed the unstructured workshop announcement, isolating 14 distinct structured entity instances with $100\%$ precision, summarized in Table~\ref{tab:exp3_regex_results}.
\begin{table}[H]
\centering
\small
\caption{Extracted Entities from Unstructured Announcement Text via Regular Expressions}
\label{tab:exp3_regex_results}
\begin{tabularx}{\textwidth}{lp{4cm}X}
\toprule
\textbf{Entity Type} & \textbf{Applied Regex Pattern} & \textbf{Extracted Text Entities} \\
\midrule
\textbf{Email Addresses} & \texttt{[a-zA-Z0-9\_.+-]+@[a-zA-Z0-9-]+\textbackslash.[a-zA-Z0-9-.]+} & \texttt{nlpworkshop@gmail.com}, \texttt{support@python.org} \\
\textbf{Web URLs} & \texttt{https?://[\textasciicircum\textbackslash s]+|www\textbackslash .[\textasciicircum\textbackslash s]+} & \texttt{https://www.nlpworkshop.com}, \texttt{www.python.org.}, \texttt{https://github.com/NLPWorkshop} \\
\textbf{Phone Numbers} & \texttt{\textbackslash+?\textbackslash d\{1,3\}[-.\textbackslash s]?\textbackslash d\{10\}|\textbackslash b\textbackslash d\{10\}\textbackslash b} & \texttt{+91-9876543210}, \texttt{9123456789} \\
\textbf{Hashtags} & \texttt{\#\textbackslash w+} & \texttt{\#NLP}, \texttt{\#Python}, \texttt{\#MachineLearning} \\
\textbf{Mentions} & \texttt{@\textbackslash w+} & \texttt{@gmail}, \texttt{@python}, \texttt{@NLPWorkshop}, \texttt{@PythonLearner} \\
\bottomrule
\end{tabularx}
\end{table}

\section{Conclusion}
Experiment~3 successfully realized all learning and technical objectives:
\begin{enumerate}
    \item Heuristic stemming via the Porter Stemmer demonstrated rapid execution but suffered from crude suffix truncation, generating non-words (\texttt{studi}) and failing completely on irregular inflections (\textit{went, mice, ate}).
    \item Morphological lemmatization via WordNet proved to be vastly superior in preserving valid linguistic lemmas, though its performance on non-nominal words is strictly contingent upon providing accurate Part-of-Speech tags.
    \item Suppletive morphological changes (\textit{better} $\to$ \texttt{good}, \textit{went} $\to$ \texttt{go}) can only be resolved by lexicon-backed lemmatizers with syntactic context.
    \item Deterministic regular expressions provided an efficient, rule-based approach for information extraction, capturing electronic communications, URLs, contact numbers, and social metadata with high fidelity.
\end{enumerate}
These insights confirm that lemmatization should be preferred over stemming in applications requiring grammatical fidelity, such as question answering and machine translation, whereas stemming suffices for search indexing where computational speed is paramount.
